\section{Fundamentação teórica}\label{sec:fundamentacao_teorica}

Esta seção apresenta algumas constatações sobre a pesquisa na área de impressões digitais e técnicas de extração de características.

\subsection{Biometria por Impressão Digital}\label{subsec:impressoes_digitais}

O termo ``Impressão Digital'' no contexto de biometria por impressões digitais se refere ao padrão de pequenos cumes e vales que podem ser observados na ponta dos dedos das pessoas. Ao longo da evolução da espécie humana, o desenvolvimento de digitais está relacionado à capacidade de segurar objetos com as mãos. O padrão específico de cumes e vales, assim como outras características do corpo, é definido por fatores genéticos e ambientais~\cite{handbook_of_fp_recognition}.

O reconhecimento de que as impressões digitais são únicas para cada dedo de cada pessoa ocorreu em 1880 com base em observações experimentais de Henry Faulds~\cite{FingerPrintTechniques}. Algum tempo depois, a unicidade de impressões digitais foi reconhecida por entidades governamentais do Reino Unido, e foi determinado que as digitais de criminosos seriam coletadas no momento da prisão. Assim, seria possível identificar pessoas a partir de impressões digitais coletadas em cenas de crime a partir da análise feita por especialistas em investigação forense. Conforme explica~\citeonline{handbook_of_fp_recognition}, apesar dos esforços para tornar mais eficiente o processo manual de identificação por impressões digitais, o aumento da demanda tornou esse método inviável. O sistema de classificação manual era ineficiente, o treinamento de novos profissionais capazes de empregá-lo era demorado e a comparação visual de impressões digitais era cansativa e lenta. Por isso, os órgãos de segurança pública começaram a investir em soluções eletrônicas e automatizadas, o que levou ao desenvolvimento dos Sistemas Automatizados de Identificação de Impressões Digitais (\textit{Automated Fingerprint Identification Systems} - AFIS).

Inicialmente usados por forças de segurança, esses sistemas hoje também são aplicados em diversas áreas não forenses devido ao aumento das preocupações com segurança e fraudes de identidade. Por exemplo, sistemas biométricos têm sido adotados em \textit{Internet Banking} e autenticação de usuários em dispositivos móveis. Nesse contexto, as digitais são coletadas com sensores especializados que geram imagens das impressões digitais, que são processadas a fim de se obter uma forma de representação que pode ser usada para verificação~\cite{handbook_of_fp_recognition}.


\subsection{Formas de representação de impressões digitais}\label{subsec:representacao_de_digitais}

\citeonline{handbook_of_fp_recognition} enfatiza que uma boa representação de uma impressão digital deve: ($i$) conter informações que possibilitem distinguir entre duas amostras; ($ii$) ser extraída de forma rápida e fácil; e ($iii$) ser compacta, de modo a poder ser armazenada em sistemas biométricos com pouca capacidade. Representações baseadas em imagens, em geral, não são adequadas devido às variações de iluminação, resolução, qualidade e aos ruídos provocados por cicatrizes ou outras lesões nos dedos, especialmente em pessoas que realizam trabalhos manuais com frequência.


\subsection{Extração de características de digitais}\label{sec:revisao_bibliografica}

Com o objetivo de selecionar os métodos que seriam avaliados neste projeto, foi feita uma revisão bibliográfica sobre artigos que descreviam métodos de extração que, preferencialmente, utilizassem algoritmos de Aprendizado de Máquina e que utilizassem conjuntos de dados públicos para treinamento e teste dos modelos, para facilitar a reprodutibilidade.

Na Tabela~\ref{tab:revisao-bibliografica} há alguns dos artigos lidos e as técnicas de extração de características propostas.

{ \footnotesize
\begin{tabularx}{\textwidth}{| l | X | X | l |}
    \caption{Revisão Bibliográfica}\label{tab:revisao-bibliografica} \\
    \midrule
    \textbf{Referência} & \textbf{Conjunto de dados} & \textbf{Métricas} & \textbf{Código aberto} \\ 
    \midrule
    \endhead
    \citeonline{mcc} & FVC2006 & $EER$, $FMR$ & Sim \\
    \midrule
    \citeonline{MinutiaeNet} & FVC2002, FVC2004, NIST SD27 & F1 Score, Precisão, \textit{Recall} & Sim \\
    \midrule
    \citeonline{deep_print} & Privado, NIST SD4, NIST SD14, FVC2004 & \textit{Acc}, $FMR$ & Sim \\
    \midrule
    \citeonline{latent-fingerprint-segmentation} & NIST SD27, WVU, IIITD & MDR, FDR, SA & Não \\
    \midrule
    \citeonline{fcn} & Privado, NIST SD27 & \textit{Recall}, Precisão, F1 Score & Não \\
    \midrule
    \citeonline{fingernet} & Privado, NIST SD27 & Precisão, \textit{Recall} & Sim \\
    \midrule
    \citeonline{highres_fingerprint_matching} & Privado, FVC2002, FVC2004 & $EER$, $FMR$, $FNMR$ & Não \\
    \midrule
    \citeonline{fcnp} & Privado, NIST SD4, NIST SD14 & Top-K ER & Não \\
    \midrule
    \citeonline{flare} & Privado, NIST SD4, NIST SD14, SD27, SD302/N2N, FVC2002, FVC2004, FVC2006 & Curvas DET, TAR com FAR fixa & Sim \\
    \bottomrule
\end{tabularx}
}

Em~\cite{mcc} foi proposta uma forma de representação chamada de~\textit{Minutia Cylinder-Code}. Ela é baseada em estruturas de dados 3D, os cilindros, que codificam os ângulos das minúcias e as distâncias entre elas.

\citeonline{MinutiaeNet} criaram a \textit{MinutiaeNet}, que é um modelo composto por duas redes neurais convolucionais, a~\textit{CoarseNet}, para estimar os mapas de \textit{scores} de minúcias com base em informações de domínio, e a~\textit{FineNet}, para refinamento dos~\textit{scores}.

\citeonline{deep_print} introduziram a~\textit{DeepPrint}, uma CNN que extrai uma representação de tamanho fixo utilizando alinhamento por Transformadores Espaciais~\cite{spatial-transformer-networks}, informações de domínio e Aprendizado Multitarefa~\cite{multitask-learning}.

\citeonline{latent-fingerprint-segmentation} criaram uma técnica para segmentação de impressões digitais utilizando uma rede neural profunda composta por pilhas de Máquinas de Boltzmann Restritas (\textit{Restricted Boltzmann Machine}, RBM) para delimitar impressões digitais latentes classificando seções das imagens entre duas classes: ``impressão digital'' e ``\textit{background}''. 

\citeonline{fcn} utilizaram uma Rede Neural Totalmente Convolucional (\textit{Fully Convolutional Network}, FCN) para selecionar seções da imagem com uma probabilidade alta de conter uma minúcia. As seções selecionadas serviram de entrada para uma CNN dedicada à classificação propriamente dita, que também calculava a orientação das minúcias.

\citeonline{fingernet} criaram a~\textit{FingerNet}, que é uma CNN que extrai um mapa de minúcias utilizando informações de domínio como orientação dos cumes e segmentação, que também são extraídas pela rede, além de realçar as digitais.

\citeonline{highres_fingerprint_matching} utilizaram duas CNNs baseadas na arquitetura ResNet~\cite{resnet-34}, uma que extrai uma representação de tamanho fixo e outra um mapa de minúcias. A verificação é feita através de cálculo de distâncias de ambas as representações.

\citeonline{fcnp} desenvolveram uma FCN que extrai uma representação de tamanho fixo com base em atributos globais e locais. O treinamento é feito utilizando porções da imagem (chamadas de ``\textit{patches}'') centradas nas minúcias. Na fase de indexação, a rede recebe a imagem completa da digital, e utiliza~\textit{Global Average Pooling} (GAP) para agregar todas as informações locais extraídas pela FCN em uma única representação global.

Em~\cite{flare} é apresentado o FLARE, um framework de reconhecimento de impressões digitais baseado na extração de representações de tamanho fixo na forma de um tensor. O método opera em três etapas integradas: inicialmente, alinha a imagem em um sistema de coordenadas padronizado utilizando métodos complementares de estimativa de pose 2D; a seguir, submete a imagem a duas estratégias de aprimoramento (UNetEnh e PriorEnh) que melhoram a clareza das cristas e suprimem o ruído de fundo preservando a textura original da modalidade da captura; por fim, uma rede neural extrai um descritor denso tridimensional que foca exclusivamente na área útil da impressão digital, codificando com alta precisão as informações espaciais globais e locais (minúcias).

Dentre as técnicas que utilizam redes neurais profundas para extração de características, as que se destacaram são o mapa de minúcias, que é uma representação de tamanho variável que armazena a posição de cada minúcia e sua orientação, e as representações de tamanho fixo, que podem ser os~\textit{embeddings} das camadas totalmente conectadas das CNNs imediatamente antes do cálculo das probabilidades de cada uma das classes (normalmente de qual dedo a digital foi coletada), isto é, um vetor unidimensional, ou um descritor denso multidimensional, capaz de agregar mais características das digitais, como segmentação e posição das minúcias. Conforme explica~\citeonline{deep_print}, a vantagem das representações de tamanho fixo é a necessidade de menos poder computacional para verificação, que pode ser feita por um simples cálculo de distância e comparação com um limiar de corte, e de espaço para armazenamento. O DeepPrint e o FLARE são exemplos de técnicas que extraem representações de tamanho fixo.

Outro aspecto que esteve presente nos experimentos foi a utilização de informações de domínio, que são as características intrínsecas das digitais~\cite{handbook_of_fp_recognition}. Quando essas informações são incluídas no treinamento dos modelos ou aprendidas em conjunto com as representações o seu desempenho tende a ser superior ao obtido com atributos mais abstratos, como apenas a textura das digitais.

Nos experimentos desenvolvidos nos artigos foram utilizados conjuntos de dados públicos e privados, com diferentes características. Por exemplo,~\citeonline{highres_fingerprint_matching} utilizou impressões digitais parciais de alta resolução, o que torna viável a inclusão de informações sobre os atributos de Nível 3 na rede neural. Já em~\cite{fcn} foram utilizadas impressões digitais latentes, que são coletadas em cenas de crimes quando pessoas tocam em objetos~\cite{handbook_of_fp_recognition} e normalmente apresentam uma alta quantidade de ruídos. No caso do FLARE foram necessários diversos bancos de imagens diferentes e técnicas de~\textit{Data Augmentation} para treinar cada uma das redes neurais que compõem o~\textit{framework}, como o~\text{Diverse Pose Fingerprint} (DPF), com ampla variação de pose, e imagens de alta resolução com degradações simuladas (variações de umidade, cicatrizes e desfoque).

Destas técnicas, foram encontradas implementações do MCC,~\textit{MinutiaeNet},~\textit{FingerNet},~\textit{DeepPrint} e FLARE. Além das implementações das técnicas, foi encontrado o modelo pré-treinado do~\textit{DeepPrint}, disponibilizado por~\citeonline{Rohwedder2023} no site~\footnote{\url{https://github.com/tim-rohwedder/fixed-length-fingerprint-extractors}}, e também dos módulos que compõem o FLARE, como os pesos das redes de realce UNetEnh e PriorEnh disponíveis no site~\footnote{\url{https://github.com/Yu-Yy/FLARE_ENH}}, as redes de alinhamento VotingPose e RegressionPose e a rede de extração da representação densa, FDD, disponíveis no site~\footnote{\url{https://github.com/Yu-Yy/FLARE}}, que são repositórios públicos do próprio autor do artigo.